\uuid{iPGb}
\titre{Changement de variable}
\auteur{hanani ; mzari}
\organisation{exo7}
\licence{CC BY-NC-SA 4.0}
\datecreate{2020-01-20}
\niveau{L1}
\module{Analyse}
\chapitre{Calcul d'intégrales}
\sousChapitre{Autre}
\difficulte{3}
\theme{Calculs d'intégrales, changement de variable}
\exo7qcmid{lille-2.integrales.015}

\contenu{
\begin{question}
\qclassification{127.99}
\qtags{motcle=changement de variable}
On note $\displaystyle I=\int _{0}^{\ln \sqrt{3}}\frac{\mathrm{e}^x\, \mathrm{d}x}{1+\mathrm{e}^{2x}}$ et $\displaystyle J=\int _{0}^{\ln 2}\frac{\mathrm{d}x}{1+\mathrm{e}^{x}}$. Le changement de variable $t=\mathrm{e}^x$ donne :
\begin{answers}
\good{$\displaystyle I=\int _1^{\sqrt{3}}\frac{\mathrm{d}t}{1+t^2}=\frac{\pi}{12}$.}
\bad{$\displaystyle I=\int _1^{\sqrt{3}}\frac{t\mathrm{d}t}{1+t^2}=\frac{1}{2}\ln 2$.}
\bad{$\displaystyle J=\int _1^{2}\frac{\mathrm{d}t}{1+t}=\ln 3-\ln 2$.}
\good{$\displaystyle J=\int _1^{2}\frac{\mathrm{d}t}{t(1+t)}=2\ln 2-\ln 3$.}
\end{answers}
\vskip3mm
\begin{explanations}
Avec $t=\mathrm{e}^x\Rightarrow \mathrm{d}t=\mathrm{e}^x\, \mathrm{d}x$, on a : $\displaystyle I=\int _1^{\sqrt{3}}\frac{\mathrm{d}t}{1+t^2}=\Big[\arctan t\Big]_1^{\sqrt{3}}=\frac{\pi}{12}$.
\vskip0mm
Avec $t=\mathrm{e}^x\Rightarrow \mathrm{d}t=\mathrm{e}^x\, \mathrm{d}x$, on a : $\displaystyle J=\int _1^{2}\frac{\mathrm{d}t}{t(1+t)}=\left[\ln \frac{t}{1+t}\right]_1^{2}=2\ln 2-\ln 3$.
\end{explanations}
\end{question}
}
